\section{Editing and cache validity}\label{sec:C12:cache}
Text editing changes the input sequence, while cache reuse assumes particular computational dependencies remain valid. In a standard causal transformer, a cached hidden state for a later token can depend on earlier tokens through attention and residual computation. Deleting an earlier span does not generally leave later states unchanged, even if the surviving text is byte-identical. Reusing cached values therefore requires a justified dependency rule, not merely matching the visible suffix.

For a simple prefix-cache implementation, an edit at position $j$ invalidates the reusable prefix from that position onward. If $L$ is the new length, the re-prefill work depends on $L-j$ as well as newly generated tokens. A context policy that repeatedly edits early content can reduce token count while destroying reuse. This explains why a representation optimized only for shortness can be expensive under an actual serving architecture. The relevant objective includes both retained size and edit location.

PagedAttention addresses logical and physical allocation and sharing of cached blocks~\cite{kwon2023pagedattention}. It does not, by itself, establish that arbitrary semantic edits preserve the numerical validity of a cached block. A page allocator may retain bytes correctly while the model computation those bytes represent is no longer appropriate for the edited sequence. Allocation validity and inference validity must be checked separately.

CLM Suffix Cache Reuse is an explicit approximation: surviving suffix states retain dependencies on the pre-edit prefix~\cite{shao2026context}. Its full-attention path re-rotates relocated keys for their new positions, while the studied hybrid path resumes pre-edit recurrent state instead of recomputing inserted content in linear-attention layers. Position adjustment does not remove stale content dependencies. Report this quality--compute tradeoff separately from exact prefix reuse and from measured end-to-end latency.

An edit-aware serving contract can attach a dependency fingerprint to every reusable region. The fingerprint identifies model version, tokenizer, position convention, attention mask, and relevant predecessor state. A reuse decision is valid only if the implementation's equivalence conditions hold. Exact fingerprints may be expensive; conservative invalidation sacrifices reuse to preserve correctness. Approximate reuse is possible, but its approximation should be evaluated as a quality tradeoff rather than described as lossless without proof.

Measure context-control latency end to end. Include the editing model call, storage writes, cache invalidation, re-prefill, extra retrieval, and any retries after a damaged summary. Report warm and cold behavior separately. A policy that looks strong when every prompt receives a fresh prefill can perform differently in a service with extensive prefix reuse. The comparison must reflect the intended deployment instead of selecting the serving regime most favorable to one method.

